\begin{proposition}[Proved under stated hypotheses: rational exceptional trees]
\label{D:rational-tree-Qfactorial}
Let $A$ be the local ring of a normal surface over an algebraically closed
field $k$ of characteristic $p>0$.  Let $f:Y\to\operatorname{Spec}A$
be a projective resolution, isomorphic away from the closed point.
Assume that its reduced exceptional divisor $E=\bigcup_{i=1}^rE_i$
has simple normal crossings, that each $E_i$ is isomorphic to
$\mathbf P^1_k$, and that its dual graph is a tree.
Then $A$ is $\mathbf Q$-factorial.  More precisely, there is an integer
$a\geq0$, depending on this resolution, such that
\[
  |\det(E_i\cdot E_j)|p^a\,\operatorname{Cl}(A)=0.
\]
The proposition does not require that $A$ have a rational singularity.
For a regular local ring the conclusion holds with annihilator one.
\end{proposition}

\begin{proof}
We assume that the exceptional divisor is nonempty, and write
$M=(E_i\cdot E_j)$.  The exceptional intersection matrix is negative
definite, with nonnegative off-diagonal entries
\cite[Section 10.1]{Kollar2013D}.
There is a positive integral exceptional cycle $Z$ such that
$Z\cdot E_i<0$ for every $i$.  Here is a direct linear-algebra argument.
Put $B=-M$ and solve $Bv=\mathbf1$ over $\mathbf Q$.
If $v=v_+-v_-$ is the decomposition into its nonnegative positive and
negative parts and $v_-\ne0$, then
\[
 0<v_-^t\mathbf1=v_-^tBv_+-v_-^tBv_-<0,
\]
because the supports of $v_+$ and $v_-$ are disjoint and the off-diagonal
entries of $B$ are nonpositive.  Thus $v\geq0$.
If $v_i=0$, then $(Bv)_i\leq0$, again a contradiction.
Multiplying $v$ by a positive common denominator gives the coefficients
of the required $Z$.

The invertible sheaf $\mathcal O_Z(-Z)$ has positive degree on every
reduced component of $Z$, and is therefore ample
\cite[Lemma 33.44.15, tag 0B5Y]{C_StacksFormalPic}.
Serre vanishing supplies an integer $N\geq1$ such that
\[
 H^1(Z,\mathcal O_Z(-nZ))=0\qquad(n\geq N)
\]
\cite[Lemma 30.17.1, tag 0B5U]{C_StacksFormalPic}.
The thickening $nZ\subset(n+1)Z$ has square-zero ideal
$\mathcal O_Z(-nZ)$.  Its units--Picard exact sequence shows that
\[
 \operatorname{Pic}((n+1)Z)\longrightarrow\operatorname{Pic}(nZ)
 \quad\hbox{is injective},
\]
and that
\[
 H^0((n+1)Z,\mathcal O_{(n+1)Z}^{\times})
   \longrightarrow H^0(nZ,\mathcal O_{nZ}^{\times})
 \quad\hbox{is surjective}
\]
for every $n\geq N$
\cite[Lemma 37.4.1, tag 0C6R]{C_StacksFormalPic}.
In particular, the first finite thickening $NZ$ detects a line bundle
which is trivial on every higher thickening; the surjectivity on units
also permits compatible choices of its trivializations.

Let $\mathcal I\subset\mathcal O_{NZ}$ be the ideal of $E$.
It is nilpotent.  Choose $a\geq0$ such that
$\mathcal I^{p^a}=0$.  The sheaf of abelian groups $1+\mathcal I$
is killed by $p^a$, since $(1+x)^{p^a}=1+x^{p^a}$ in characteristic $p$.
The units sequence
\[
 1\longrightarrow1+\mathcal I\longrightarrow
 \mathcal O_{NZ}^{\times}\longrightarrow\mathcal O_E^{\times}
 \longrightarrow1
\]
therefore shows that $p^a$ kills
$\ker[\operatorname{Pic}(NZ)\to\operatorname{Pic}(E)]$.

We claim that every invertible sheaf $\mathcal L$ on $Y$ having
degree zero on all $E_i$ satisfies $\mathcal L^{\otimes p^a}\simeq
\mathcal O_Y$.  A degree-zero invertible sheaf on $E$ is trivial:
it is trivial on each $\mathbf P^1$ component, and its gluing constants
at the nodes can be removed successively along the tree.
Thus $\mathcal L^{\otimes p^a}|_{NZ}$ is trivial.  The injectivity
above shows that its restriction to every $nZ$ is trivial, and
surjectivity on units makes these trivializations compatible for
$n\geq N$.

First suppose that $A$ is complete.  The powers of the ideal of $Z$
and the powers of $\mathfrak m_A\mathcal O_Y$ define the same adic
topology, because their radicals agree on the Noetherian scheme $Y$.
The theorem on formal functions in degree zero thus algebraizes the
compatible trivializing sections to a section of
$\mathcal L^{\otimes p^a}$ on $Y$
\cite[Theorem 30.20.5, tag 02OC]{C_StacksFormalPic}.
Its vanishing locus is closed and disjoint from $E$.
Its proper image in the local scheme $\operatorname{Spec}A$ is
therefore empty.  Hence the section is everywhere nonvanishing.

For general $A$ we apply this argument after the faithfully flat
base change $A\to\widehat A$.  The cycles $nZ$ and all calculations
on them are unchanged, since each is annihilated by some power of
$\mathfrak m_A$.  Put $\mathcal M=\mathcal L^{\otimes p^a}$.
Properness and flat base change give
\[
 H^0(Y,\mathcal M)\otimes_A\widehat A
   =H^0(Y_{\widehat A},\mathcal M_{\widehat A})\simeq\widehat A,
\]
where we also use $f_*\mathcal O_Y=\mathcal O_{\operatorname{Spec}A}$
\cite[Lemma 30.5.2, tag 02KH]{C_StacksFormalPic}.
Consequently the finite $A$-module $H^0(Y,\mathcal M)$ is free of
rank one.  Its evaluation map on $Y$ is an isomorphism after the same
faithfully flat base change, and hence is an isomorphism before base
change.  This proves the claim.

Finally let $D$ be a Weil divisor on $\operatorname{Spec}A$, and
let $\widetilde D$ be its strict transform.  Since $Y$ is regular,
$\widetilde D$ is Cartier.  Inverting the matrix $M$ supplies rational
numbers $b_i$ such that
\[
 \left(\widetilde D+\sum_i b_iE_i\right)\cdot E_j=0
 \qquad(1\leq j\leq r).
\]
For $m=|\det M|$ all $mb_i$ are integers.  The Cartier divisor
$m\widetilde D+\sum_i mb_iE_i$ thus defines an invertible sheaf
which is killed by $p^a$, by the claim.  Its $p^a$ multiple is a
principal divisor on $Y$.  Pushing that principal divisor down to
$\operatorname{Spec}A$ shows that $mp^aD$ is principal.  This proves
the asserted bounded exponent and the $\mathbf Q$-factorial conclusion.
\end{proof}

\begin{remark}[Source scope]
For one irreducible exceptional curve, the numerical correction and
descent mechanism also occur in Tanaka's Propositions 3.25 and 3.26
\cite[pp. 23--24]{C_Tanaka2014}.  Those propositions are stated for
one exceptional curve.  The argument above handles the entire rational
tree at once and explicitly accounts for the nilpotent thickenings.
It makes no assertion that a rational exceptional tree forces
$R^1f_*\mathcal O_Y=0$.
\end{remark}
